% !TeX root = ../main.tex
\section{Электродинамическое проектирование полосового фильтра на диэлектрическом резонаторе (HFSS)}

Для реализации частотной селекции сигналов в выходном тракте усилителя на центральной частоте $f_0 = 4{,}5\text{ ГГц}$ в среде \textit{Ansys HFSS} была разработана 3D электродинамическая модель фильтра на диэлектрическом цилиндрическом резонаторе (ДР), связанном с микрополосковой линией передачи (рис.~\ref{fig:dr_geometry}).

\begin{figure}[H]
    \centering
    \includegraphics[width=0.75\textwidth]{05_sch.png}
    \caption{3D-геометрия полосового фильтра на диэлектрическом резонаторе в Ansys HFSS}
    \label{fig:dr_geometry}
\end{figure}

\subsection{Электродинамический пересчет геометрии резонатора}

Исходная базовая модель резонатора была рассчитана на центральную частоту $f_{\mathrm{orig}} = 3{,}5\text{ ГГц}$ и обладала следующими геометрическими параметрами:
\begin{itemize}
    \item Радиус диэлектрического цилиндра: $R_{\mathrm{orig}} = 5{,}2\text{ мм}$;
    \item Высота цилиндра: $H_{\mathrm{orig}} = 10{,}4\text{ мм}$;
    \item Относительная диэлектрическая проницаемость материала резонатора: $\varepsilon_r = 10$.
\end{itemize}

Поскольку для диэлектрического резонатора на основной рабочей моде колебаний резонансная частота обратно пропорциональна его линейным размерам ($f_0 \propto \frac{1}{D \sqrt{\varepsilon_r}}$), для пересчета резонатора на целевую частоту усилителя $f_0 = 4{,}5\text{ ГГц}$ был применен масштабный коэффициент геометрического подобия:
\begin{equation}
    k = \frac{f_{\mathrm{orig}}}{f_0} = \frac{3{,}5\text{ ГГц}}{4{,}5\text{ ГГц}} \approx 0{,}778.
\end{equation}

Скорректированные размеры цилиндра резонатора составили:
\begin{itemize}
    \item Радиус цилиндра: $R = R_{\mathrm{orig}} \cdot k = 5{,}2\text{ мм} \cdot 0{,}778 \approx 4{,}05\text{ мм}$;
    \item Высота цилиндра: $H = H_{\mathrm{orig}} \cdot k = 10{,}4\text{ мм} \cdot 0{,}778 \approx 8{,}10\text{ мм}$.
\end{itemize}

\subsection{Результаты моделирования S-параметров фильтра}

Электродинамический расчет матрицы рассеяния фильтра выполнен методом конечных элементов (FEM) со следующими параметрами адаптивного сетапа:
\begin{itemize}
    \item Частота адаптации сетки (\textit{Solution Frequency}): $4{,}5\text{ ГГц}$;
    \item Тип частотного свипа: \textit{Interpolating Sweep};
    \item Диапазон частот: от $3{,}5\text{ ГГц}$ до $5{,}5\text{ ГГц}$ (401 расчетная точка).
\end{itemize}

Полученные частотные зависимости коэффициентов отражения $|S_{11}|$ и передачи $|S_{21}|$ представлены на рис.~\ref{fig:dr_sparam}.

\begin{figure}[H]
    \centering
    \includegraphics[width=0.85\textwidth]{04_plot.png}
    \caption{Частотные зависимости параметров $S_{11}$ и $S_{21}$ спроектированного фильтра на ДР}
    \label{fig:dr_sparam}
\end{figure}

По результатам электродинамического анализа зафиксировано:
\begin{itemize}
    \item Коэффициент отражения $|S_{11}|$ демонстрирует глубокий резонансный минимум на частоте $4{,}45\text{ ГГц}$, достигающий значения $-22{,}5\text{ дБ}$, что свидетельствует об эффективном согласовании тракта на резонансной частоте;
    \item Коэффициент передачи $|S_{21}|$ формирует полосно-пропускающий пик с максимумом в полосе частот $4{,}5\dots 4{,}6\text{ ГГц}$ и уровнем затухания в полосе порядка $-7\text{ дБ}$;
    \item При отстройке от центральной частоты более чем на $1\text{ ГГц}$ затухание фильтра составляет более $20\dots 25\text{ дБ}$, что обеспечивает требуемую селективность для последующей интеграции в тракт усилителя.
\end{itemize}
